\section{Discussion}
\label{sec:discussion}

\paragraph{Where the agent was reliable, and where not}
It was reliable at everything that ended in an artefact a machine could check: the Lean
development has no gaps because Lean would not accept them, the tables match the result files
because a validator reads both, the negative results survived because the directory structure made
deleting them harder than keeping them. It was unreliable at everything that ended in a sentence.
Incident~1 is the clearest: a manuscript that passed $162\,327$ checks and was, in the researcher's
words, poor, because the prose had become a transcript of checks the agent had itself proposed. Incident~4 is the subtler one: the agent saw an anomaly, produced an
explanation of the right shape, and accepted it. Both are failures of judgment, invisible to every
mechanism except a person reading.

\paragraph{Design consequences}
Three follow. First, every claim should be routed toward an artefact check: if a constant is
asserted, write the experiment that measures it (Incident~7); if two descriptions of the same object
exist, require them to agree (the theorem). Second, validators should check facts, not prose: the
frozen ledger pinned phrases and shaped the paper; its replacement pins numbers to files. Third, the
human's attention should be spent on outputs, not process: the human-caught incidents came from
reading a page and looking at a figure, and none of the \nHuman{} messages needed expertise beyond
recognizing that something looked wrong.
